%%
%% 03_literature_review.tex — Literature Review
%%

\section{Literature Review}\label{sec:literature_review}

\subsection{Sarcasm Detection}\label{subsec:sarcasm_detection}

Early sarcasm detection methods used feature engineering (punctuation, word patterns, sentiment flip) and traditional machine learning (SVM, Random Forest). Deep learning and hybrid models later dominated the field. For example, Sharma~\emph{et~al.}~\cite{sharma2020deep} trained CNNs and LSTMs on tweets and Reddit posts, achieving approximately 90\% accuracy. Multimodal approaches incorporate images alongside text (e.g.\ a ``comic'' dialogue context). Transformer-based models (BERT~\cite{devlin2019bert}, RoBERTa~\cite{liu2019roberta}, XLNet, etc.) have more recently set new state-of-the-art results. The parent work~\cite{oprea2025llm} finds that fine-tuned transformers achieve macro F1 scores in $[0.878\text{--}0.943]$ on sarcasm corpora. Particularly, DistilBERT~\cite{sanh2019distilbert} pre-trained on sentiment (SST-2) gave the most stable $0.8784$ macro-F1 across diverse domains.

\subsection{Transformer Models and Baselines}\label{subsec:transformer_baselines}

Comparative studies show deep transformers outperform LSTMs and feature-based methods in sarcasm detection. For example, RoBERTa-large and RoBERTa-base fine-tuned on headlines achieved 94.3\% accuracy ($\text{F1} = 0.9425$), whereas smaller models like DistilBERT trade off some performance for efficiency. Other works have evaluated BERT with BiLSTM, ensemble CNN-LSTM, etc., on various languages (Hindi, Arabic) and domains. In summary, transformer fine-tuning on in-domain data is now standard practice, motivating us to reuse exactly the setup of~\cite{oprea2025llm} for compatibility.

\subsection{Explainability in NLP and Sarcasm}\label{subsec:xai_nlp}

\textbf{Post-hoc explanation techniques} have become common in NLP to interpret black-box models. LIME~\cite{ribeiro2016why} and SHAP~\cite{lundberg2017unified} generate token importance scores via perturbations or Shapley values. Integrated Gradients~\cite{sundararajan2017axiomatic} and attention visualization are also used. In sarcasm detection specifically, Kumar~\emph{et~al.}~\cite{kumar2021explainable} applied XGBoost with LIME/SHAP on the MUStARD dialogue corpus; their explanations highlight the words driving the model's sarcastic/not judgment. Bagate~\emph{et~al.}~\cite{bagate2025sarcasm} used an LSTM+SVC fusion on a self-annotated Reddit political dataset, and implemented counterfactual explanation: by permuting words, they identify which word change flips the model's sarcasm prediction. They find counterfactual highlights intuitive cues (e.g.\ the word whose removal breaks the sarcasm). Both works report that explanation (words highlighted by LIME/SHAP or counterfactual) aids human understanding, but neither integrates generative rationales.

Attention-based interpretability has been explored as well: some models deliberately enforce interpretable attention (e.g.\ multi-head self-attention with sparsity) for sarcasm. However, attention weights alone may not yield human-plausible explanations (they often require aggregation or thresholding). Recent surveys of NLP XAI~\cite{bilal2025llms} note the diversity of explanation formats: token importance, gradients, or \emph{natural language rationales}. Indeed, increasingly researchers are using LLMs themselves to \textbf{generate textual explanations}.

Bueno~\emph{et~al.}~\cite{bueno2026scoring} propose combining model-agnostic feature attribution (SHAP) with LLM-generated rationales, and find that while LLM explanations are fluent, SHAP attributions are typically more faithful to model decisions. Inoue~\emph{et~al.}~\cite{inoue2026lime} introduce LIME-LLM, using an LLM (Anthropic Claude) to produce fluent counterfactual edits for LIME; their method outperforms vanilla LIME/SHAP and matches integrated gradients in aligning with human rationales. Meanwhile, Li~\emph{et~al.}~\cite{li2025drift} show that naively prompting LLMs yields unfaithful reasoning, and propose a dual-reward ``Drift'' method to improve faithfulness of LLM rationales to model behaviour. Our work is inspired by these trends: we adopt LLM prompting to generate explanations post-hoc, but we also critically evaluate faithfulness and compare to established XAI methods.

\subsection{Literature Comparison}\label{subsec:lit_comparison}

Table~\ref{tab:literature_comparison} contrasts key recent studies (including the parent paper) on sarcasm detection and explainability.

\begin{sidewaystable}[p]
\caption{Comparison of recent studies on sarcasm detection and explainability.}

\label{tab:literature_comparison}
\centering
\footnotesize
\renewcommand{\arraystretch}{2}
\newcolumntype{L}[1]{>{\RaggedRight\hyphenpenalty=10000\exhyphenpenalty=10000\arraybackslash}p{#1}}
\begin{tabular}{L{1.8cm} L{2.4cm} L{2.2cm} L{2.2cm} L{2.6cm} L{2.6cm} L{2.6cm}}
\toprule
\textbf{Authors (Year)} &
\textbf{Dataset} &
\textbf{Model} &
\textbf{Explainability} &
\textbf{Advantages} &
\textbf{Limitations} &
\textbf{Contributions} \\
\midrule
Oprea \& B\^{a}ra (2025) &
News Headlines (Kaggle, 55k), Amazon Reviews (Sarcasm Amazon) &
RoBERTa-large, RoBERTa-base, DistilBERT, DistilBERT-SST2 (fine-tuned) &
\textit{None (baseline classification)} &
\textbf{Systematic evaluation} of multiple transformers; group-aware splits; high macro-F1 ($>$0.94). &
No model interpretability beyond a reproducible classification pipeline. &
Provides a robust and reproducible baseline for sarcasm classification. \\
\midrule
Kumar et al. (2021) &
MUStARD dialogue corpus &
XGBoost using ensemble features &
LIME and SHAP (post-hoc token importance) &
\textbf{Context-aware dialogue} sarcasm detection with word-level explanations using LIME and SHAP. &
Relies on handcrafted features and classical machine learning rather than deep neural models. &
Demonstrates that contextual information improves sarcasm detection while providing interpretable explanations. \\
\midrule
Bagate et al. (2025) &
Self-annotated Reddit political corpus &
LSTM + SVC fusion &
Counterfactual explanation using word permutations &
\textbf{Domain-specific} political sarcasm detection; ensemble improves F1; counterfactuals identify influential words. &
Counterfactual generation is computationally expensive and explanations remain word-based. &
Introduces counterfactual explainability for identifying words responsible for sarcasm. \\
\midrule
Bueno et al. (2026) &
Teaching transcripts (CLASS dataset) &
Fine-tuned transformer PLMs and prompted LLMs &
SHAP and LLM rationale evaluation &
\textbf{Compares SHAP with LLM-generated explanations} and proposes a systematic evaluation framework. &
Focuses on educational scoring rather than sarcasm; LLM explanations found less faithful. &
Introduces deletion-based evaluation methods for assessing explanation faithfulness. \\
\midrule
Inoue et al. (2026) &
SST-2, CoLA, HateXplain &
LLMs (Claude, GPT) &
LIME-LLM (LLM-generated counterfactuals) &
\textbf{Generative counterfactual explanations} outperform conventional LIME/SHAP and align well with human rationales. &
Requires repeated LLM inference for every instance and has not been evaluated on sarcasm datasets. &
Shows that LLM-generated counterfactuals substantially improve explanation fidelity. \\
\midrule
\textit{This work (2026)} &
Same datasets as Oprea \& B\^{a}ra (2025) &
Same fine-tuned transformers &
LLM-generated textual rationales (post-hoc) &
Produces \textbf{human-readable explanations} of sarcasm cues while maintaining classification performance. &
Dependent on LLM availability; explanation faithfulness must be quantitatively validated. &
Extends the existing sarcasm detection pipeline with an LLM-based explainability module that bridges prediction and interpretation. \\
\bottomrule
\end{tabular}
\end{sidewaystable}

This survey shows a clear evolution: earlier sarcasm detection focused on accuracy~\cite{oprea2025llm} or basic interpretability~\cite{kumar2021explainable,bagate2025sarcasm}. Recent trends incorporate \textbf{LLM-based explanations}~\cite{bueno2026scoring,inoue2026lime} and evaluate faithfulness. Our proposed work builds on these advances by applying free-form rationale generation specifically for sarcasm, guided by known cues (contextual incongruity, sentiment flip, irony markers) in a domain-specific yet generalizable way.
